\documentclass[preview]{standalone}

\usepackage{amsmath}
\usepackage{amssymb}
\usepackage{stellar}
\usepackage{definitions}

\begin{document}

\id{unique-factorization-domains}
\genpage

\section{Factorization elements}

\begin{snippetdefinition}{domain-factorization-elements-definition}{Associates, irreducible elements, and prime elements}
    Let \(A\) be an \integraldomain.
    \begin{enumerate}
        \item Elements \(a,b\in A\) are \emph{associates} if
        \(b=au\) for some unit \(u\in A^\times\). In this case one writes
        \(a\sim b\).
        \item An element \(a\in A\) is \emph{irreducible} if
        \(a\neq0_A\), \(a\notin A^\times\), and every factorization
        \(a=xy\) has \(x\in A^\times\) or \(y\in A^\times\).
        \item An element \(a\in A\) is \emph{prime} if
        \(a\neq0_A\), \(a\notin A^\times\), and the \ideal \((a)\) is
        prime. Equivalently,
        \[
            a\divides xy
            \implies
            a\divides x\ \lor\ a\divides y
        \]
        for all \(x,y\in A\).
    \end{enumerate}
\end{snippetdefinition}

\begin{snippet}{irreducible-prime-element-distinction}
    An \irreducibleelement has no nontrivial factorization. A
    \primeelement satisfies the stronger divisibility property that, when
    it divides a product, it divides at least one factor.
\end{snippet}

\begin{snippetproposition}{associate-irreducible-prime-properties}{Associates, irreducibles, and primes}
    Let \(A\) be an \integraldomain.
    \begin{enumerate}
        \item For \(a,b\in A\),
        \[
            a\sim b
            \ifandonlyif
            (a)=(b).
        \]
        \item Every \primeelement is \irreducibleelement[irreducible].
        \item If \(A\) is a \principalidealdomain, every
        \irreducibleelement is prime.
    \end{enumerate}
\end{snippetproposition}

\begin{snippetproof}{associate-irreducible-prime-properties-proof}{associate-irreducible-prime-properties}{Associates, irreducibles, and primes}
    If \(b=au\) for a unit \(u\), then \(b\in(a)\) and
    \(a=bu^{-1}\in(b)\), so \((a)=(b)\). Conversely, suppose
    \((a)=(b)\). If \(a=0_A\), then \(b=0_A\), and the two elements are
    associates. Otherwise write
    \[
        b=ac,
        \qquad
        a=bc'.
    \]
    Then \(a=acc'\), so cancellation gives \(cc'=1_A\). Thus \(c\) is a
    unit and \(a,b\) are \associateelements.

    Let \(a\) be prime and suppose \(a=xy\). Since \(xy\in(a)\),
    primality gives \(x\in(a)\) or \(y\in(a)\). If \(x=ac\), then
    \[
        a=xy=acy,
    \]
    so cancellation gives \(cy=1_A\), and \(y\) is a unit. The other case
    similarly makes \(x\) a unit. Hence \(a\) is irreducible.

    Finally, let \(A\) be a PID and let \(a\) be irreducible. If
    \[
        (a)\subseteq J\subseteq A,
    \]
    write \(J=(b)\). Then \(a=bc\) for some \(c\in A\). Irreducibility
    implies that \(b\) or \(c\) is a unit. In the first case \(J=A\); in
    the second, \(a\) and \(b\) are associates, so \(J=(a)\). Thus
    \((a)\) is maximal, hence prime, and \(a\) is a \primeelement.
\end{snippetproof}

\section{Unique factorization}

\begin{snippetdefinition}{unique-factorization-domain-definition}{Unique factorization domain}
    A \emph{unique factorization domain}, abbreviated \emph{UFD}, is an
    \integraldomain \(A\) in which every nonzero nonunit can be written as
    a finite product of \primeelement[prime elements]. More precisely, if
    \[
        a=p_1\cdots p_r=q_1\cdots q_s
    \]
    are two such factorizations, then \(r=s\) and, after reordering,
    \(p_i\sim q_i\) for every \(i\).
\end{snippetdefinition}

\begin{snippetproposition}{prime-factorization-uniqueness}{Uniqueness of factorization into prime elements}
    In an \integraldomain, a factorization into
    \primeelement[prime elements] is unique up to reordering and replacing
    factors by \associateelements.
\end{snippetproposition}

\begin{snippetproof}{prime-factorization-uniqueness-proof}{prime-factorization-uniqueness}{Uniqueness of factorization into prime elements}
    Suppose
    \[
        p_1\cdots p_r=q_1\cdots q_s,
    \]
    with every factor prime. Since \(p_1\) divides the product on the
    right, it divides some \(q_j\). After reordering, assume
    \(q_1=p_1u\). The element \(q_1\) is irreducible and \(p_1\) is not a
    unit, so \(u\) is a unit. Hence \(q_1\sim p_1\), so they are
    \associateelements.

    Substitution and cancellation of the nonzero element \(p_1\) give
    \[
        p_2\cdots p_r=uq_2\cdots q_s.
    \]
    Absorb the unit \(u\) into one of the remaining prime factors and
    repeat. This process pairs all factors with \associateelements. Since no prime
    factor is a unit, neither side can have unpaired factors; therefore
    \(r=s\).
\end{snippetproof}

\begin{snippettheorem}{principal-ideal-domain-is-unique-factorization-domain}{Every PID is a UFD}
    Every \principalidealdomain is a \uniquefactorizationdomain.
\end{snippettheorem}

\begin{snippetproof}{principal-ideal-domain-is-unique-factorization-domain-proof}{principal-ideal-domain-is-unique-factorization-domain}{Every PID is a UFD}
    Let \(A\) be a PID. Every \ideal of \(A\) is finitely generated, so
    \(A\) is a \noetherianring. Suppose that some nonzero nonunit is not a
    finite product of primes, and let \(\mathcal{S}\) be the nonempty
    \set of principal \ideal[ideals] generated by such elements. The
    maximal-element characterization of Noetherian rings gives a maximal
    member \((a)\in\mathcal{S}\).

    The element \(a\) is not prime. In a PID every irreducible is prime,
    so \(a\) is reducible:
    \[
        a=bc
    \]
    for nonunits \(b,c\). Then
    \[
        (a)\subseteq(b),
        \qquad
        (a)\subseteq(c).
    \]
    Both inclusions are strict. For example, equality
    \((a)=(b)\) would make \(a\) and \(b\) associates; comparing
    \(a=bc\) with \(a=bu\) and cancelling \(b\neq0_A\) would make \(c\)
    a unit.

    Maximality of \((a)\) in \(\mathcal{S}\) therefore implies that both
    \(b\) and \(c\) are finite products of primes. Hence so is
    \(a=bc\), a contradiction. Thus every nonzero nonunit factors into
    primes, and uniqueness follows from the preceding proposition.
\end{snippetproof}

\begin{snippetcorollary}{euclidean-domain-is-unique-factorization-domain}{Every Euclidean domain is a UFD}
    Every \euclideandomain is a \uniquefactorizationdomain.
\end{snippetcorollary}

\begin{snippetproof}{euclidean-domain-is-unique-factorization-domain-proof}{euclidean-domain-is-unique-factorization-domain}{Every Euclidean domain is a UFD}
    Every \euclideandomain is a \principalidealdomain, and every
    \principalidealdomain is a \uniquefactorizationdomain.
\end{snippetproof}

\end{document}
